\documentclass{article}
\usepackage{graphicx}
\usepackage[margin=1in]{geometry}
\usepackage{amsmath, amssymb, amsthm}
\usepackage{mathtools}
\usepackage{bm}
\usepackage{hyperref}


\DeclareMathOperator*{\clip}{clip}
\newcommand{\Ehat}{\widehat{\mathbb{E}}}


\title{2026 NIPS}
\author{YAREN ZHANG}
\date{April 2026}

\begin{document}
\maketitle

\section{Multi Domain}

The document below presents the multi-domain BSPO objective in a bottom-up
order: (1) notation, (2) bisimulation-distance estimates, (3) policy-ratio
statistics $s^{\pm}$, (4) clipping operators $\Delta^{(\cdot)}$, (5) threshold
hyperparameters $\delta$, (6) penalty hyperparameters $\lambda$, (7) advantage
estimators $\hat A$, and finally (8) the five objective variants together with
(9) an optional token-level refinement.

% =============================================================================
\section{Notation and hierarchy}

The policy is optimised over a three-level hierarchy.

\begin{align*}
\text{Hierarchy indices:}\quad & \text{D}, \; \operatorname{Grp}, \; \operatorname{Tj} \\
\text{Position indices:}\quad & dom, \; g, \; \tau, \; t \\
\text{Generic position index across hierarchy:}\quad & z
\end{align*}

Throughout, $\Ehat_{t\in\tau\sim g\sim dom}[\cdot]$ denotes the empirical
expectation over tokens inside trajectories, trajectories inside groups, and
groups inside domains (hierarchical nesting read outside-in).

% =============================================================================
\section{Bisimulation-distance estimates}

Both the $\delta$-thresholds (\S\ref{sec:delta}) and the Wasserstein form of
the advantage (\S\ref{sec:adv}) depend on a plug-in estimate of a
bisimulation-style distance at each hierarchy level.

\begin{align*}
\hat d(z) &= \mathbb{E}_{\tau \sim \mathcal{P}(z)}\bigl[\,|R(\tau) - \mathbb{E}_{\tau' \sim \mathcal{P}(z)}[R]|\,\bigr] \\[3pt]
R_{\max\text{res}}(dom) &= \max_{\tau \in dom}\,|R(\tau) - \bar R_{dom}|
\end{align*}

Group-level bisimulation size is estimated from the relative group/domain
token counts:
\[
\hat d_{\operatorname{Grp}}(g)
  = \frac{{T_g}^{2}}{{T_{dom}}^{2}}\,\hat d(g).
\]

% =============================================================================
\section{Policy-ratio statistics \texorpdfstring{$s^{\pm}$}{s+/-}}

Let $r_t = \pi_\theta(a_t\mid s_t) / \pi_{\theta_{\text{old}}}(a_t\mid s_t)$.
One-sided deviations are aggregated upward through the hierarchy as an
\emph{empirical expectation at every level}, matching the mean-per-level
reduction pattern of verl's \verb|seq-mean-token-mean| aggregator:

\begin{align*}
s_\tau^{\pm} \equiv s_{\operatorname{Tj}}^{\pm}(\tau)
  &= \frac{1}{|\tau|}\sum_{t\in\tau}\max\!\bigl(\pm(r_t-1),\,0\bigr),
     \qquad s_\tau = s_\tau^{+} - s_\tau^{-} \\[3pt]
s_g^{\pm}   &= \frac{1}{|g|}\sum_{\tau \in g}\max(\pm s_\tau,\,0) \\[3pt]
s_{dom}^{\pm} &= \frac{1}{|dom|}\sum_{g \in dom}\max(\pm s_g,\,0)
\end{align*}

where $|g|$ is the number of trajectories in group $g$ (16 in the reference
Phase-1 ablation) and $|dom|$ is the number of groups in domain $dom$. At
every level, $s_z = s_z^{+} - s_z^{-}$.

\paragraph{Why mean at every level.}
Keeping all three statistics on the same O(1) scale has three payoffs.
(i) A single $\delta$ scale (the per-token arithmetic deviation inherited
from GSPO) transfers up the hierarchy without per-level cascade factors;
this also simplifies the $\delta_z$ derivation of \S\ref{sec:delta}.
(ii) The reduction aligns with verl's infrastructure: $s_g^{\pm}$ and
$s_{dom}^{\pm}$ become plain \verb|scatter_mean| calls keyed on the
prompt-index / data-source tensors that verl already carries in the
non-tensor batch.
(iii) The telescoping identity used in Variant~2 (\S\ref{sec:obj}) requires
that $\widehat{\mathbb E}_{\tau\sim g\sim dom}[X(dom)] = \widehat{\mathbb
E}_{dom}[X]$, which holds only when the aggregation at each level weights
members uniformly.

% =============================================================================
\section{Clipping operators \texorpdfstring{$\Delta_z^{(\cdot)}$}{Delta}}

Two clipping styles are used by different objectives:

\begin{align*}
\Delta_z^{(\text{W})}   &= \operatorname{sgn}(s_z)\,\clip\!\bigl(\min(s_z^{+},s_z^{-}),\,\delta_z\bigr) \\
\Delta_z^{(\text{reg})} &= \clip(s_z^{+},\delta_z) - \clip(s_z^{-},\delta_z)
\end{align*}

\subsection*{Why \texorpdfstring{$\Delta_z^{(\text{W})}$}{Delta\^{}W} keeps the $\min(s_z^+, s_z^-)$ magnitude, and how $\operatorname{sgn}(s_z)$ is handled}

This subsection explains (i) why the magnitude must be
$\min(s_z^+, s_z^-)$ rather than $|s_z|$ or $s_z$, (ii) what $\operatorname{sgn}(s_z)$ is doing
mathematically, and (iii) how the torch implementation handles the
$\operatorname{sgn}$ non-differentiability without changing the math.

\paragraph{(a) The magnitude is the Wasserstein coupling, $\min(s_z^+,s_z^-)$.}
Interpret the per-token log-ratios $\{r_t - 1\}_{t\in\tau}$ as a signed
measure on the real line: $s_\tau^+ = \tfrac{1}{|\tau|}\sum_t[r_t-1]_+$
is the total ``upward'' mass and $s_\tau^- = \tfrac{1}{|\tau|}\sum_t[1-r_t]_+$
is the total ``downward'' mass. The optimal transport coupling between
the old-policy ratio distribution (a Dirac at $1$) and the new-policy
ratio distribution (spread over $\{r_t\}$) splits into two contributions:
\[
\underbrace{\min(s_\tau^+,\,s_\tau^-)}_{\text{symmetric transport}}
\;+\;
\underbrace{|s_\tau|}_{\text{net translation}},
\]
which together sum to the total transport cost $s_\tau^+ + s_\tau^- =
\tfrac{1}{|\tau|}\sum_t|r_t-1|$. The first term, $\min(s_\tau^+,s_\tau^-)$,
counts mass that is moved both up and down in equal measure
(``bidirectional churn''); the second, $|s_\tau|$, counts net drift. In
the BSPO formulation, $\Delta_z^{(\text{W})}$ isolates the \emph{symmetric}
portion, because that is the contribution that cannot be absorbed into
the net-drift term $s_z$ and must be controlled by $\delta_z$ as a
hard Wasserstein budget.

Concretely, consider a trajectory where half the tokens have
$r_t=1.1$ and half have $r_t=0.9$: $s_\tau^+=s_\tau^-=0.05$,
$s_\tau=0$. The policy has drifted substantially at every token but the
net per-trajectory drift is zero. The Wasserstein bisimulation distance
is $0.05$ (the symmetric-transport cost); the magnitude
$\min(s_\tau^+,s_\tau^-) = 0.05$ reports it. Replacing the magnitude with
$|s_\tau| = 0$ or with $s_\tau = 0$ would lose this signal and allow the
policy to drift arbitrarily in the balanced-two-sided regime without
triggering the clip or penalty.

\emph{Conclusion.} The magnitude of $\Delta_z^{(\text{W})}$ must be
$\min(s_z^+,s_z^-)$ (clipped at $\delta_z$) to preserve the Wasserstein
coupling interpretation. It is not interchangeable with $|s_z|$.

\paragraph{(b) The $\operatorname{sgn}(s_z)$ factor is a direction selector, not a learnable quantity.}
The magnitude $\clip(\min(s_z^+,s_z^-),\delta_z)$ is non-negative. A
sign must be attached to align $\Delta_z^{(\text{W})} \hat A_z$ with
the policy-gradient convention: a trajectory with net upward drift
($s_z>0$) and positive advantage should contribute positively to the
objective; same direction, negative advantage should contribute
negatively. $\operatorname{sgn}(s_z)$ encodes ``direction of net drift''
--- a \emph{discrete} categorical fact about the batch, determined
entirely by whether more tokens went up or down. It is not a continuous
parameter that the optimiser should learn to move.

\paragraph{(c) The sgn non-differentiability is an implementation issue, not a math flaw.}
$\operatorname{sgn}:\mathbb R\to\{-1,0,+1\}$ has derivative zero
a.e.\ (piecewise constant) and is undefined at $s_z=0$. In autograd this
means the factor $\operatorname{sgn}(s_z)$ contributes no gradient on any
path. This is the mathematically correct behaviour: we do not want to
learn ``move towards the other sign.'' The magnitude
$\clip(\min(s_z^+,s_z^-),\delta_z)$ carries all the gradient signal
about the bisimulation distance.

The only concern is that autograd may emit an undefined / NaN value at
exactly $s_z=0$ if the backward is naively computed through
$\operatorname{sgn}$. We sidestep this by explicitly telling autograd to
treat the sign as a constant per-sample:
\begin{verbatim}
sign_s = torch.sign(s_z).detach()
delta_w = sign_s * torch.clamp(
    torch.minimum(s_pos, s_neg), max=delta_z
)
\end{verbatim}
The \verb|.detach()| is not a math modification; it is a declaration
that the sign is a \emph{categorical direction}, and instructs autograd
to stop the backward chain at the sign node (gradient $=0$, which is
the mathematical a.e.\ derivative of $\operatorname{sgn}$ anyway).

\paragraph{(d) Why not replace $\operatorname{sgn}(s_z)\cdot\clip(\min,\delta_z)$ with $\clip(s_z,-m_z,m_z)$.}
One might try to sidestep the sign by writing
$\Delta_z^{(\text{W})} = \clip(s_z,-m_z,+m_z)$ with
$m_z = \clip(\min(s_z^+,s_z^-),\delta_z)$. The two forms agree whenever
$|s_z|\geq m_z$ (clipping regime). When $|s_z|<m_z$ they differ:
$\operatorname{sgn}\cdot\clip(\min,\delta_z)$ returns $\pm m_z$ (full
symmetric-transport magnitude), while $\clip(s_z,-m_z,m_z)$ returns
$s_z$ (net drift). The latter is \textbf{wrong} by the argument of
paragraph (a): it discards the Wasserstein coupling precisely in the
balanced-two-sided regime where that coupling is the entire
bisimulation signal. The example of half-up, half-down trajectories
above produces $\Delta=\pm 0.05$ under the correct form and
$\Delta=0$ under the replacement; the correct form preserves the
constraint, the replacement removes it. The replacement is therefore
ruled out on theoretical grounds.

\paragraph{(e) If a fully-differentiable sign surrogate is required.}
For settings where \verb|.detach()| is not available or a smooth
backward is desired, replace $\operatorname{sgn}(s_z)$ with
\[
\operatorname{sgn}_\varepsilon(s_z) \;:=\; \frac{s_z}{\sqrt{s_z^2+\varepsilon^2}},
\qquad \varepsilon\ll\delta_z,
\]
which is smooth, bounded in $(-1,+1)$, agrees with $\operatorname{sgn}$
whenever $|s_z|\gg\varepsilon$, and has subgradient
$\varepsilon^2/(s_z^2+\varepsilon^2)^{3/2}$ everywhere. The cost is an
additional hyperparameter $\varepsilon$ and a spike of gradient through
the sign path near $s_z=0$; the detach approach is preferred as default
because the mathematical derivative of $\operatorname{sgn}$ is zero
a.e.\ and detaching reports exactly that.

\paragraph{(f) Worked examples --- value of $\Delta_z^{(\text{W})}$ across regimes.}
Fix $\delta_z = 0.1$. For each $(s_z^+, s_z^-)$ we compute
$s_z = s_z^+-s_z^-$, $\min(s_z^+,s_z^-)$, $m_z = \clip(\min,\delta_z)$,
and the correct output $\Delta_z^{(\text{W})} = \operatorname{sgn}(s_z)\cdot m_z$.

\begin{center}
\renewcommand{\arraystretch}{1.35}
\begin{tabular}{c | c c c c c | c | l}
 & $s_z^+$ & $s_z^-$ & $s_z$ & $\min$ & $m_z$ & $\Delta_z^{(\text{W})}$ & interpretation \\ \hline
A & $0.08$ & $0.02$ & $+0.06$ & $0.02$ & $0.02$ & $+0.02$ & net-up dominated; $s^-$ sets the symmetric floor \\
B & $0.08$ & $0.05$ & $+0.03$ & $0.05$ & $0.05$ & $+0.05$ & net-up small, but bidirectional churn is $0.05$ \\
C & $0.02$ & $0.08$ & $-0.06$ & $0.02$ & $0.02$ & $-0.02$ & mirror of A; $s^+$ is the smaller side \\
D & $0.05$ & $0.08$ & $-0.03$ & $0.05$ & $0.05$ & $-0.05$ & mirror of B; Wasserstein says $0.05$ matters \\
E & $0.10$ & $0.10$ & $0.00$  & $0.10$ & $0.10$ & $0\pm$ & fully balanced; net drift undefined in direction \\
F & $0.30$ & $0.20$ & $+0.10$ & $0.20$ & $0.10$ & $+0.10$ & $\delta_z$ cap binds; full budget spent up \\
G & $0.30$ & $0.25$ & $+0.05$ & $0.25$ & $0.10$ & $+0.10$ & $\delta_z$ cap binds; symmetric churn saturated \\
\end{tabular}
\end{center}

Notes on the rows:
\begin{itemize}
\item Rows A and C are the ``standard'' operating regime (one side
dominates, magnitude determined by the smaller side which serves as the
bisimulation floor).
\item Rows B and D are the \emph{critical} regime where BSPO departs
from plain $s_z$-based clipping: net drift is small ($0.03$) but the
symmetric churn is larger ($0.05$). Wasserstein bisimulation assigns
the larger value to the trajectory, and $\Delta_z^{(\text{W})}$ reports
it. Replacing the magnitude with $|s_z|=0.03$ or $s_z=\pm 0.03$ would
under-constrain the policy by $40\%$ in this regime.
\item Row E is the balanced on-policy state. The direction is
ill-defined (by convention $\operatorname{sgn}(0)=0$, so
$\Delta_z^{(\text{W})} = 0$); any infinitesimal perturbation flips
$\Delta_z^{(\text{W})}$ to $\pm 0.10$, reflecting the discrete nature
of the direction selector. The \verb|.detach()| convention treats this
as the a.e.\ derivative, gradient $=0$ at $s_z=0$.
\item Rows F and G show the $\delta_z$-cap binding. In row F, net drift
and $\delta_z$-capped magnitude coincide. In row G, net drift ($0.05$)
is smaller than the $\delta_z$ cap but the symmetric churn ($0.25$) was
already saturated by the cap, so $\Delta_z^{(\text{W})} = +\delta_z = +0.10$.
\end{itemize}

\paragraph{(g) Gradient flow under the correct form.}
With $\operatorname{sgn}(s_z)$ detached, the backward chain is
\[
\frac{\partial (\Delta_z^{(\text{W})} \hat A_z)}{\partial \theta}
\;=\; \operatorname{sgn}(s_z)\,\hat A_z \cdot \frac{\partial\,\clip(\min(s_z^+,s_z^-),\,\delta_z)}{\partial \theta},
\]
where $\partial\,\clip(\cdot,\delta_z)/\partial\theta = 0$ when the
$\delta_z$ cap binds, and otherwise $\partial\min(s_z^+,s_z^-)/\partial\theta$
routes gradient to whichever of $s_z^+, s_z^-$ is the smaller side
(via \verb|torch.minimum| subgradient), then through the
ratio $\to\log\pi_\theta$ chain.

Only the magnitude path carries gradient; the sign is a read-only
feature of the current batch. This is the correct behaviour both
mathematically (sgn has zero derivative a.e.) and semantically
(direction of net drift is not a learnable degree of freedom).

\paragraph{(h) Summary.}
$\Delta_z^{(\text{W})} = \operatorname{sgn}(s_z)\cdot\clip(\min(s_z^+,s_z^-),\delta_z)$
is the correct mathematical form: the magnitude $\min(s_z^+,s_z^-)$
encodes the Wasserstein symmetric-transport contribution and cannot be
replaced by $|s_z|$ or $s_z$ without losing the bisimulation signal in
balanced-drift regimes. The $\operatorname{sgn}$ factor is a discrete
direction indicator whose mathematical derivative is zero a.e.; its
autograd handling via \verb|torch.sign(s).detach()| realises this zero
derivative explicitly and avoids undefined behaviour at $s_z=0$. No
smooth surrogate is required for the default implementation.

Expanded at each hierarchy level, the regularised form is:

\begin{align*}
\Delta_\tau^{(\text{reg})}
  &= \clip\!\bigl(s_{\operatorname{Tj}}^{+}(\tau),\,\delta_\tau\bigr) - \clip\!\bigl(s_{\operatorname{Tj}}^{-}(\tau),\,\delta_\tau\bigr) \\
\Delta_g^{(\text{reg})}
  &= \clip\!\bigl(s_{\operatorname{Grp}}^{+}(g),\,\delta_g\bigr) - \clip\!\bigl(s_{\operatorname{Grp}}^{-}(g),\,\delta_g\bigr) \\
\Delta_{dom}^{(\text{reg})}
  &= \clip\!\bigl(s_d^{+}(dom),\,\delta_{dom}\bigr) - \clip\!\bigl(s_d^{-}(dom),\,\delta_{dom}\bigr)
\end{align*}

% =============================================================================
\section{Threshold hyperparameters \texorpdfstring{$\delta$}{delta}}
\label{sec:delta}

Three flat (domain-invariant) hyperparameters are provided to the training
script: $\delta_{\text{D}}$, $\delta_{\operatorname{Grp}}$,
$\delta_{\operatorname{Tj}}$. The token-level threshold is inherited from
GRPO. From these, the per-level effective thresholds are derived:

\begin{align*}
\delta_{dom} &= \frac{\hat d(dom)}{R_{\max\text{res}}(dom)}\,\delta_{\text{D}} \\[3pt]
\delta_g
  &= \frac{\hat d(dom)\,\delta_{\text{D}}}{R_{\max\text{res}}(dom)}
     \left( \frac{\hat d_{\operatorname{Grp}}(g)}{\sum_{g':dom} \hat d_{\operatorname{Grp}}(g')^{2}} \right)
     \delta_{\operatorname{Grp}} \\[3pt]
\delta_\tau
  &= \frac{\hat d(dom)\,\delta_{\text{D}}}{R_{\max\text{res}}(dom)}
     \left( \frac{\hat d_{\operatorname{Grp}}(g)}{\sum_{g':dom} \hat d_{\operatorname{Grp}}(g')^{2}} \right)
     \frac{\delta_{\operatorname{Tj}}}{T_g}
\end{align*}

Hyperparameter vector:
$\bigl[\delta_{\text{D}},\; \delta_{\operatorname{Grp}},\; \delta_{\operatorname{Tj}}\bigr]$,
with the token threshold from GRPO.

% =============================================================================
\section{Penalty hyperparameters \texorpdfstring{$\lambda$}{lambda}}

Only the Wasserstein-penalty objectives (\S\ref{sec:obj}, variants 4--5) use
a one-sided soft penalty whenever the min-side ratio exceeds its threshold.
The penalty weights are flat across domain/group/trajectory:
\[
\bigl[\lambda_{\text{D}},\; \lambda_{\operatorname{Grp}},\; \lambda_{\operatorname{Tj}}\bigr],
\qquad \text{default }\approx [\,1,\; 0.001,\; ?\,].
\]

% =============================================================================
\section{Advantage estimators \texorpdfstring{$\hat A$}{A-hat}}
\label{sec:adv}

The base trajectory advantage uses either the standard GRPO normalisation or,
if the Wasserstein mode is active, a bisimulation-scaled variant:
\[
\hat A_\tau
  = \frac{R_\tau - \bar R_g}{\sigma_g}
  \quad\text{or (Wasserstein)}\quad
  \hat A_\tau
  = \frac{1}{\min\!\bigl(\hat d(dom),\,1/|g(\tau)|\bigr)}\cdot\frac{R_\tau - \bar R_g}{R_{\max\text{res}}(dom)}.
\]

Hierarchical averages flow upward with $T_g / T_{dom}$ token-count weighting:
\begin{align*}
\hat A_g   &= \frac{1}{T_g}\sum_{\tau:g}\hat A_\tau \\
\hat A_{dom} &= \frac{T_g}{T_{dom}}\sum_{g:dom}\hat A_g
\end{align*}

% =============================================================================
\section{Objectives}
\label{sec:obj}

The five variants below are ordered from simplest to most general. The
symbol $128\times 16$ in the first name refers to the group / response grid
used in the reference Phase-1 ablation and is not part of the mathematical
definition.

\subsection*{1. Simplest baseline}
\begin{align*}
\hat J_{128\times 16}
  &= \Ehat_{t \in \tau \sim g \sim dom}\bigl[\Delta_\tau^{(\text{reg})}\hat A_\tau\bigr] \\
  &= \Ehat_{t \in \tau \sim g \sim dom}\Bigl[\bigl(\clip(s_z^{+},\delta_z) - \clip(s_z^{-},\delta_z)\bigr)\hat A_\tau\Bigr]
\end{align*}

\subsection*{2. Baseline with hierarchy (telescoping, regularised)}
\begin{align*}
\hat J
  &= \Ehat_{t \in \tau \sim g \sim dom}\Bigl[
        \Delta_{dom}^{(\text{reg})}\hat A_{dom}
      + \bigl(\Delta_g^{(\text{reg})}   - \Delta_{dom}^{(\text{reg})}\bigr)\hat A_g
      + \bigl(\Delta_\tau^{(\text{reg})} - \Delta_g^{(\text{reg})}\bigr)\hat A_\tau
     \Bigr] \\
  &= \Ehat_{t \in \tau \sim g \sim dom}\bigl[\bigl(\Delta_\tau^{(\text{reg})} - \Delta_g^{(\text{reg})}\bigr)\hat A_\tau\bigr]
\end{align*}

\subsection*{3. Strict Wasserstein (telescoping, min-side clip)}
\begin{align*}
\hat J
  &= \Ehat_{t \in \tau \sim g \sim dom}\Bigl[
        \operatorname{sgn}(s_{dom})\clip\!\bigl(\min(s_{dom}^{+},s_{dom}^{-}),\delta_{dom}\bigr)\hat A_{dom} \\
  &\qquad + \operatorname{sgn}(s_g)\clip\!\bigl(\min(s_g^{+},s_g^{-}),\delta_g\bigr)\hat A_g \\
  &\qquad + \operatorname{sgn}(s_\tau)\clip\!\bigl(\min(s_\tau^{+},s_\tau^{-}),\delta_\tau\bigr)\hat A_\tau
     \Bigr] \\
  &\equiv \Ehat_{t \in \tau \sim g \sim dom}\bigl[
        \Delta_{dom}^{(\text{W})}\hat A_{dom}
      + \Delta_g^{(\text{W})}\hat A_g
      + \Delta_\tau^{(\text{W})}\hat A_\tau
     \bigr] \\
  &= \Ehat_{t \in \tau \sim g \sim dom}\bigl[\Delta_\tau^{(\text{W})}\hat A_\tau\bigr]
\end{align*}

\subsection*{4. Wasserstein penalty (trajectory clip + three-level penalty)}
\begin{align*}
\hat J
  &= \Ehat_{t \in \tau \sim g \sim dom}\Bigl[
        \Delta_\tau^{(\text{W})}\hat A_\tau
      - \lambda_{\operatorname{Tj}}\max\!\Bigl(0,\tfrac{\min(s_\tau^{+},s_\tau^{-})}{\delta_\tau}-1\Bigr) \\
  &\qquad - \lambda_{\operatorname{Grp}}\max\!\Bigl(0,\tfrac{\min(s_g^{+},s_g^{-})}{\delta_g}-1\Bigr)
      - \lambda_{\text{D}}\max\!\Bigl(0,\tfrac{\min(s_{dom}^{+},s_{dom}^{-})}{\delta_{dom}}-1\Bigr)
     \Bigr]
\end{align*}

\subsection*{5. Wasserstein penalty only (raw trajectory ratio + three-level penalty)}
\begin{align*}
\hat J
  &= \Ehat_{t \in \tau \sim g \sim dom}\Bigl[
        s_\tau \hat A_\tau
      - \lambda_{\operatorname{Tj}}\max\!\Bigl(0,\tfrac{\min(s_\tau^{+},s_\tau^{-})}{\delta_\tau}-1\Bigr) \\
  &\qquad - \lambda_{\operatorname{Grp}}\max\!\Bigl(0,\tfrac{\min(s_g^{+},s_g^{-})}{\delta_g}-1\Bigr)
      - \lambda_{\text{D}}\max\!\Bigl(0,\tfrac{\min(s_{dom}^{+},s_{dom}^{-})}{\delta_{dom}}-1\Bigr)
     \Bigr]
\end{align*}

% =============================================================================
\section{Optional refinement: token-level separation at the trajectory term}

Where the trajectory objective uses the raw $s_\tau\hat A_\tau$ product
(variant 5, and the unclipped components of variants 3--4), the scalar
$s_\tau$ may be replaced by its mean plus a token-clipped residual:
\[
  s_\tau\,\hat A_\tau
  \;\Longrightarrow\;
  \left(\,s_\tau + \frac{1}{|\tau|}\sum_{t\in\tau}\clip(\Delta_t - s_\tau,\,\cdot)\right)\hat A_\tau.
\]
The second argument of $\clip$ (the per-token bound) is left unset here; a
natural choice is the token-level threshold inherited from GRPO.

\section{Single Domain}

ratios deviation:
\begin{align*}
{}^{(\text{GSPO})}\tilde{\delta}_\tau 
&= \frac{\hat d(dom)\,\delta_{\text{D}}}{R_{\max res}(dom)}
   \left( \frac{\hat d_{\operatorname{Grp}}(g)}{\sum_{g':dom} \hat d_{\operatorname{Grp}}(g')^2} \right)
   \frac{\delta_{\operatorname{Tj}}}{T_g} \\
&\doteq \frac{1}{\sum_{g':dom} \hat d_{16}(g')^2}\,
   \hat d_{128\times 16}(dom)\,\hat d_{16}(g)\,\delta_{\operatorname{sing}} \\
\hat d(z) &= \mathbb{E}_{\tau \sim \mathcal{P}(z)}\bigl[\,|R(\tau) - \mathbb{E}_{\tau'\sim \mathcal{P}(z)}[R]|\,\bigr] \\
z &= \tau \text{ in the following:} \\
\Delta_z^{(\text{W})}   &= \operatorname{sgn}(s_z)\,\operatorname{clip}(\min(s_z^+, s_z^-), \delta_z) \\
\Delta_z^{(\text{reg})} &= \operatorname{clip}(s_z^+, \delta_z) - \operatorname{clip}(s_z^-, \delta_z) \\
s_\tau &= s_\tau^+ - s_\tau^- \\
s_\tau^+ &:= \frac{1}{|\tau|}\sum_{t\in\tau}\max(r_t - 1, 0), \qquad
s_\tau^- := \frac{1}{|\tau|}\sum_{t\in\tau}\max(1 - r_t, 0)
\end{align*}

Objectives
\begin{align*}
\text{Simplest baseline: }\hat J_{128\times 16}
  &= \widehat{\mathbb{E}}_{\tau\sim g\sim dom}\bigl[\Delta_\tau^{(\text{reg})}\hat A_\tau\bigr] \\
  &= \widehat{\mathbb{E}}_{\tau\sim g\sim dom}\bigl[(\operatorname{clip}(s_z^+,\delta_z) - \operatorname{clip}(s_z^-,\delta_z))\hat A_\tau\bigr] \\
\text{Strict Wasserstein: }\hat J
  &= \widehat{\mathbb{E}}_{\tau\sim g\sim dom}\bigl[\operatorname{sgn}(s_{dom})\operatorname{clip}(\min(s_{dom}^+,s_{dom}^-),\delta_{dom})\hat A_{dom} \\
  &\qquad + \operatorname{sgn}(s_g)\operatorname{clip}(\min(s_g^+,s_g^-),\delta_g)\hat A_g \\
  &\qquad + \operatorname{sgn}(s_\tau)\operatorname{clip}(\min(s_\tau^+,s_\tau^-),\delta_\tau)\hat A_\tau\bigr] \\
  &\equiv \widehat{\mathbb{E}}_{\tau\sim g\sim dom}\bigl[\Delta_{dom}^{(\text{W})}\hat A_{dom} + \Delta_g^{(\text{W})}\hat A_g + \Delta_\tau^{(\text{W})}\hat A_\tau\bigr] \\
\text{Wasserstein Penalty: }\hat J
  &= \widehat{\mathbb{E}}_{\tau\sim g\sim dom}\Bigl[\Delta_\tau^{(\text{W})}\hat A_\tau - \lambda_{\operatorname{Tj}}\max\Bigl(0,\tfrac{\min(s_\tau^+,s_\tau^-)}{\delta_\tau} - 1\Bigr)\Bigr] \\
\text{Wasserstein Penalty only: }\hat J
  &= \widehat{\mathbb{E}}_{\tau\sim g\sim dom}\Bigl[s_\tau \hat A_\tau - \lambda_{\operatorname{Tj}}\max\Bigl(0,\tfrac{\min(s_\tau^+,s_\tau^-)}{\delta_\tau} - 1\Bigr)\Bigr]
\end{align*}

Advantages
\begin{align*}
\lambda_{\operatorname{Tj}} &= 0.001 \text{ or } 1 \\
\hat A_{dom} &= \tfrac{16}{128}\sum_{g:dom}\hat A_g, \qquad
\hat A_g = \tfrac{1}{16}\sum_{\tau:g}\hat A_\tau, \qquad
\hat A_\tau = \frac{R_\tau - \bar R_g}{\sigma_g}
\end{align*}

\end{document}
